%==============================================================================
%  CAPÍTULO — DETECCIÓN TEMPRANA DE LA CRISIS Y MECANISMOS PREVENTIVOS
%==============================================================================
\chapter{La detección temprana de la crisis y los mecanismos preventivos}
\label{cap:deteccion-temprana}

\fichaCapitulo{Lo que ocurre ---o debería ocurrir--- antes del concurso.}{Los deberes de los
administradores ante la crisis, las señales de alerta, la fase preconcursal en el derecho comparado
y su ausencia en el sistema chileno.}

%==============================================================================
\bloqueTeorico
%==============================================================================

\subsection{El problema del momento: por qué se llega tarde al concurso}

El derecho concursal, tal como se lo suele enseñar, comienza cuando la crisis ya estalló. Pero la
decisión más importante ---y la que más valor conserva o destruye--- se toma \emph{antes}: en el
momento en que el deudor reconoce la crisis y decide actuar. Un concurso abierto a tiempo, sobre una
empresa que aún conserva giro, clientes y personal, admite reorganización; el mismo concurso abierto
seis meses después, cuando la caja se agotó y los proveedores huyeron, sólo admite liquidación.

La experiencia ---y las cifras chilenas, con su abrumador predominio de liquidaciones sobre
reorganizaciones (capítulo \ref{cap:renegociacion})--- muestran que \destacar{se llega tarde}. Las
razones son conocidas: el sesgo de continuación de administradores y socios que apuestan por la
resurrección con dinero ajeno (capítulo \ref{cap:teoria-economica}); el estigma que dispara la sola
apertura pública del procedimiento; y la ausencia de incentivos ---y a veces de deberes--- para
actuar antes. La detección temprana es el intento de corregir ese retraso sistemático.

\begin{foco}[El valor se conserva antes, no durante]
Ningún procedimiento concursal, por bien diseñado que esté, crea valor: en el mejor caso lo
preserva. El valor de la empresa se conserva o se pierde en la \destacar{ventana previa} a la
insolvencia irreversible. Por eso el derecho comparado ha desplazado su atención de la
administración de la crisis a su \emph{prevención}, y por eso la mayor carencia del sistema chileno
no está en cómo tramita el concurso, sino en cuán tarde permite ---o incentiva--- iniciarlo.
\end{foco}

\subsection{Los deberes de los administradores ante la crisis}

En la sociedad en crisis, la posición de los administradores se transforma. Mientras la sociedad es
solvente, su deber fiduciario mira al interés social ---a los socios---. Cuando la insolvencia se
aproxima, el interés relevante se desplaza hacia los \destacar{acreedores}, que pasan a ser los
titulares económicos residuales del patrimonio: son ellos quienes soportarán las pérdidas de las
decisiones que se tomen.

De ese desplazamiento derivan, en los ordenamientos que lo regulan, deberes específicos: el deber de
\destacar{monitorear} la situación patrimonial, el de \destacar{no agravar} la insolvencia con nuevo
endeudamiento o apuestas riesgosas, y ---en su expresión más intensa--- el deber de
\destacar{solicitar oportunamente} la apertura del procedimiento. El derecho alemán sanciona el
retraso mediante la \emph{Insolvenzverschleppung} (dilación de la insolvencia), que puede generar
responsabilidad personal de los administradores por el daño causado a los acreedores.

\begin{alerta}[En Chile el deber de solicitar a tiempo es difuso]
La \LIR{} no impone a los administradores un deber temporalmente preciso de solicitar el concurso, ni
un régimen claro de responsabilidad por su retraso. Los controles existen ---la posposición de
personas relacionadas, las acciones revocatorias del período sospechoso, los delitos concursales del
\art{463} y siguientes del \cpen{} (capítulo \ref{cap:laboral-penal})---, pero operan
\emph{ex post}, sancionando lo ya ocurrido, no incentivando la actuación temprana. Es una diferencia
de fondo con los modelos que anudan responsabilidad al retraso.
\end{alerta}

%==============================================================================
\bloqueNormativo
%==============================================================================

\subsection{Las señales de alerta: qué mirar y cuándo}

La detección temprana descansa sobre un conjunto de \destacar{indicadores} que anticipan la crisis
antes de que se manifieste como cesación de pagos. Sistematizarlos es útil tanto para el asesor del
deudor como para el acreedor que vigila su exposición:

\begin{description}[leftmargin=0pt,style=nextline,font=\sffamily\bfseries]
  \item[Señales de liquidez] Estrechamiento del capital de trabajo, uso del máximo de las líneas de
  crédito, retraso creciente en los pagos a proveedores, refinanciamientos sucesivos de la misma
  deuda, cheques a fecha como medio ordinario de pago.
  \item[Señales patrimoniales] Pérdidas recurrentes que erosionan el patrimonio, activos
  sobrevalorados en balance, pasivos contingentes no provisionados, patrimonio neto negativo.
  \item[Señales operativas] Caída sostenida de ventas o de márgenes, pérdida de clientes o
  proveedores clave, rotación de personal directivo, conflictos societarios.
  \item[Señales de comportamiento] Postergación sistemática de cotizaciones previsionales e
  impuestos ---indicador especialmente grave, porque revela que la empresa se financia con dinero
  ajeno de rango preferente---, opacidad contable, operaciones inusuales con personas relacionadas.
\end{description}

\begin{practica}[El impago previsional como la alarma más elocuente]
De todas las señales, el \destacar{atraso en cotizaciones e impuestos de retención} es la más
reveladora, y por una razón dogmática: esos fondos no son de la empresa ---los recaudó de terceros
(capítulo \ref{cap:teoria-economica})---. Una empresa que deja de enterar cotizaciones está
financiando su operación con dinero preferente ajeno, señal casi inequívoca de que la crisis de
liquidez es estructural. Para el asesor, es el indicador que obliga a poner sobre la mesa la
alternativa concursal; para el acreedor laboral o previsional, la alarma para verificar y actuar.
\end{practica}

\subsection{Los mecanismos preventivos en el derecho comparado}

Los ordenamientos que toman en serio la prevención ofrecen instrumentos que operan \emph{antes} del
concurso, y que conviene conocer porque marcan el horizonte de la reforma chilena (capítulos
\ref{cap:comparado-espana} y \ref{cap:comparado-europa-continental}):

\paragraph{Los procedimientos de alerta.} El derecho francés fue pionero con sus
\emph{procédures d'alerte}: mecanismos que obligan a distintos actores ---el auditor
(\emph{commissaire aux comptes}), el comité de empresa, incluso el presidente del tribunal--- a dar
la voz de alarma cuando detectan hechos que comprometen la continuidad de la explotación,
provocando una deliberación de los órganos sociales antes de que la crisis se consume.

\paragraph{La negociación preventiva y confidencial.} El \emph{mandat ad hoc} y la
\emph{conciliation} franceses, o los acuerdos preconcursales españoles, permiten renegociar con los
principales acreedores \destacar{sin publicidad}, evitando el estigma que precipita la crisis. Es la
herramienta cuya ausencia se hace sentir con más fuerza en Chile, donde el procedimiento es público
desde su primer día.

\paragraph{La apertura anticipada por insolvencia inminente.} El derecho alemán habilita la apertura
del procedimiento ante la \emph{drohende Zahlungsunfähigkeit} ---insolvencia inminente, no
actual---, y la Directiva europea construye sobre el umbral de la
\destacar{probabilidad de insolvencia} sus marcos de reestructuración preventiva. La lógica común:
premiar a quien actúa antes, no exigir que la crisis se consume para poder intervenir.

\begin{alerta}[La deuda pendiente del sistema chileno]
La \LIR{} carece de procedimientos de alerta, de un instrumento de negociación preventiva
confidencial y de una apertura por insolvencia inminente claramente incentivada. La reorganización
está teóricamente disponible de forma temprana ---no exige cesación de pagos---, pero nada empuja a
usarla a tiempo ni protege la negociación previa del estigma. La reforma más transformadora del
derecho concursal chileno no está en el procedimiento, sino en lo que debería ocurrir \emph{antes}
de él.
\end{alerta}

%==============================================================================
\bloquePractico
%==============================================================================

\subsection{El diagnóstico temprano en la asesoría}

Mientras el ordenamiento no ofrezca instrumentos preventivos, el asesor debe construir la prevención
con las herramientas del derecho común y con método. La secuencia:

\begin{checklist}[Diagnóstico temprano de la empresa en dificultades]
  \item \textbf{Revisar los indicadores de alerta} ---liquidez, patrimonio, operación,
  comportamiento--- y datar el inicio real de la crisis.
  \item \textbf{Cuantificar el pasivo preferente oculto}: cotizaciones e impuestos impagos, que
  suelen ser mayores y de peor rango de lo que el deudor admite.
  \item \textbf{Evaluar la viabilidad económica} distinguiéndola de la solvencia: ¿la operación
  genera márgenes positivos, o el negocio ya no es viable? (capítulo \ref{cap:teoria-economica}).
  \item \textbf{Explorar la negociación privada} con los acreedores relevantes antes de todo
  procedimiento público, y documentarla.
  \item \textbf{Si procede el concurso, elegir la vía y el momento}: reorganización mientras haya
  empresa que reestructurar; no esperar a que sólo quede liquidar.
  \item \textbf{Documentar la diligencia de los administradores}: las decisiones adoptadas en la
  crisis serán escrutadas después por la vía revocatoria y penal.
\end{checklist}

\begin{foco}[El mejor concurso es el que se prepara con tiempo]
La calidad de un procedimiento concursal se juega en su preparación, y ésta empieza mucho antes de
la presentación. El deudor que llega al concurso con la crisis datada, el pasivo preferente
cuantificado, la viabilidad evaluada y la negociación previa documentada tiene opciones; el que llega
empujado por un embargo, sólo tiene la liquidación. La detección temprana no es una fase del
procedimiento: es la condición de que el procedimiento sirva de algo. En palabras del análisis
comparado de \textcite{ruz-derecho-concursal}, el eje del derecho concursal moderno se ha desplazado
de la sanción del fracaso a la gestión anticipada de la crisis.
\end{foco}
